\documentclass[11pt]{article}
\usepackage[utf8]{inputenc}
\usepackage[T1]{fontenc}
\usepackage[english]{babel}
\usepackage{lmodern}
\usepackage{amsmath,amssymb,amsthm,mathtools}
\usepackage{geometry,booktabs,enumitem}
\usepackage[most]{tcolorbox}
\usepackage[colorlinks=true,linkcolor=blue!50!black,urlcolor=blue!50!black,citecolor=black]{hyperref}
\geometry{margin=2.6cm}
\setlist{itemsep=2pt,topsep=4pt}
\setlength{\parskip}{3pt}
\setlength{\emergencystretch}{2em}

\theoremstyle{plain}
\newtheorem{proposition}{Proposition}[section]
\newtheorem{lemme}[proposition]{Lemma}
\theoremstyle{definition}
\newtheorem{definition}[proposition]{Definition}
\newtheorem{exemple}[proposition]{Example}
\newtheorem{exercice}[proposition]{Exercise}
\theoremstyle{remark}
\newtheorem*{remarque}{Remark}

\newcommand{\JJ}{\mathcal{J}}
\newcommand{\DD}{\mathcal{D}}
\newcommand{\E}{\mathbb{E}}
\newcommand{\R}{\mathbb{R}}
\newcommand{\N}{\mathcal{N}}
\newcommand{\tr}{\operatorname{tr}}
\newcommand{\KL}[2]{D_{KL}\!\left(#1\,\Vert\,#2\right)}
% Numbers from the replication package (scripts/jfec_build.py)
\newcommand{\BIASTWENTY}{0.39}
\newcommand{\BIASFORTYEIGHT}{2.41}
\newcommand{\BIASTWENTYMC}{0.39}
\newcommand{\BIASTWENTYT}{0.39}
\newcommand{\BIASTWENTYLW}{0.002}
\newcommand{\FESHRINK}{61\%}
\newcommand{\MCGARCHIID}{17\%}
\newcommand{\MCGARCHFHS}{6.5\%}
\newcommand{\MCPOWFHS}{68\%}


\newtcolorbox{intuition}[1][]{colback=blue!3,colframe=blue!40!black,
  fonttitle=\bfseries,title={In plain words},boxrule=0.5pt,arc=2pt,#1}
\newtcolorbox{papier}[1][]{colback=orange!4,colframe=orange!60!black,
  fonttitle=\bfseries,title={Where this is used in the paper},boxrule=0.5pt,arc=2pt,#1}
\newtcolorbox{piege}[1][]{colback=red!3,colframe=red!50!black,
  fonttitle=\bfseries,title={Common pitfall},boxrule=0.5pt,arc=2pt,#1}

\title{The mathematics, step by step\\[4pt]
\large Companion note to ``What Survives a Placebo? Entropy, Market Coupling
and the Plausibility of Stress Scenarios''\\[2pt]
\normalsize (version prepared for the \emph{Journal of Financial Econometrics})}
\author{Sébastien Gallet}
\date{November 2026}

\begin{document}
\maketitle

\begin{abstract}
\noindent This note works through all the mathematics of the journal version
of the paper and of its online appendix, one step at a time. Compared with the
note on the central-bank version, the section on inference is new: the bias of
sample coupling (Bartlett's formula), why resampling $p$-values are exact, the
volatility-preserving GARCH null, the within-year gap, the Monte Carlo check
of size and power, Holm's multiple-testing procedure and the Clark--West test. It is written for a master's
student who knows probability (expectation, density, multivariate normal) and
linear algebra (determinant, trace, eigenvalues, positive definite matrices),
but not necessarily information theory. Each result is stated, proved,
illustrated with a numerical example, and linked to the exact place in the
paper where it is used. Blue boxes give the intuition, orange boxes the link
with the paper, red boxes the mistakes to avoid. Most sections end with a few
solved exercises.
\end{abstract}

\tableofcontents

%======================================================================
\section*{The storyline on one page}
\addcontentsline{toc}{section}{The storyline on one page}

The paper rests on a chain of seven ideas. Each gets its own section in this
note.

\begin{enumerate}
\item \textbf{Entropy measures the uncertainty} of a distribution, in nats
(section~\ref{sec:entropie}).
\item \textbf{The Kullback--Leibler divergence measures the gap between two
distributions}, also in nats, and it does not depend on units
(section~\ref{sec:kl}).
\item \textbf{For a fixed mean and covariance, the Gaussian is the
maximum-entropy distribution}: the ``most ignorant'' one. Any departure from
the Gaussian is therefore information (section~\ref{sec:maxent}).
\item \textbf{The entropy of a return vector splits exactly into three
pieces}: scale (volatilities), dependence (correlations) and shape (tails and
skewness). Dropping scale gives the structural index $\JJ$, which is itself a
divergence (sections~\ref{sec:gauss} to~\ref{sec:canaux}).
\item \textbf{Receiving information lowers uncertainty on average}
(a theorem), and we postulate that it then rebuilds slowly: hence a sawtooth
dynamic (section~\ref{sec:dynamique}; in the journal version these
postulates are in the online appendix).
\item \textbf{But the estimator produces this sawtooth by itself.} Placebo
nulls are therefore needed: one without dynamics, one that keeps volatility
dynamics but not correlation dynamics. Only the coupling of markets under
stress survives both (section~\ref{sec:temoin}).
\item \textbf{The nat becomes a unit of scenario severity}: $k$ standard
deviations cost $k^2/2$ nats, any scenario can be priced and the cheapest one
found (section~\ref{sec:scenarios}); and the coupling measure splits exactly
by market segment, which gives the entropic CISS (section~\ref{sec:blocs}).
\end{enumerate}

\paragraph{Notation.} $\log$ is the natural logarithm. $X\in\R^N$ is a vector
of returns with density $p$, mean $\mu$ and covariance $\Sigma$. We write
$\Sigma=DRD$ with $D=\operatorname{diag}(\sigma_1,\dots,\sigma_N)$ the
volatilities and $R$ the correlation matrix. $\N(\mu,\Sigma)$ is the normal
distribution; $q$ always denotes a Gaussian reference.

%======================================================================
\section{Entropy and its unit, the nat}
\label{sec:entropie}

\subsection{The discrete case: counting questions}

For a discrete variable taking values $x_i$ with probabilities $p_i$, the
\emph{Shannon entropy} is
\[
H(p)=-\sum_i p_i\log p_i .
\]
A few values worth knowing: a fair coin has $H=\log 2\approx0.693$; a uniform
distribution over $M$ values has $H=\log M$; a certain variable has $H=0$.
Entropy is largest for the uniform distribution and zero for a distribution
concentrated on one point.

\begin{intuition}
$H$ measures ``how much we do not know''. $-\log p_i$ is the \emph{surprise}
of outcome $i$: a rare outcome is very surprising. Entropy is the average
surprise.
\end{intuition}

\paragraph{The unit.} Choosing the base of the logarithm \emph{is} choosing
the unit. With $\log_2$ we count in bits; with the natural logarithm, in
\emph{nats}. Since $\log_2 x=\log x/\log 2$, $1~\text{nat}=1/\log2\approx1.443$
bits. A bit is the information that halves the set of possibilities; a nat
divides it by $e\approx2.718$. The whole paper is in nats.

\subsection{The continuous case: differential entropy}

For a density $p$ on $\R^N$, define
\[
h(p)=-\int p(x)\log p(x)\,dx = -\E_p[\log p(X)] .
\]

\begin{exemple}[Uniform and normal distributions in one dimension]
\label{ex:unif-norm}
The method is the same in both cases: write $\log p(x)$, then take the
expectation of $-\log p(X)$.

\medskip
\noindent\textbf{Uniform on $[0,a]$.}
\begin{enumerate}[label=\arabic*.]
\item The density is $p(x)=1/a$ on $[0,a]$ and $0$ elsewhere.
\item On $[0,a]$, $\log p(x)=\log(1/a)=-\log a$: a constant.
\item The expectation of a constant is the constant itself, so
$h=-\E[\log p(X)]=-(-\log a)=\log a$.
\end{enumerate}
Reading: the wider the interval, the larger the uncertainty. For $a=2$,
$h=\log2\approx0.69$; for $a=1$, $h=0$; for $a=1/2$, $h=-\log2<0$ (more on
this below).

\medskip
\noindent\textbf{Normal $\N(\mu,\sigma^2)$.}
\begin{enumerate}[label=\arabic*.]
\item The density is
$p(x)=\dfrac{1}{\sqrt{2\pi\sigma^2}}\exp\!\Bigl(-\dfrac{(x-\mu)^2}{2\sigma^2}\Bigr)$.
\item Take the logarithm. The log of a product is the sum of the logs, and
$\log\exp(y)=y$:
\[
\log p(x)=\log\frac{1}{\sqrt{2\pi\sigma^2}}-\frac{(x-\mu)^2}{2\sigma^2}
=-\tfrac12\log(2\pi\sigma^2)-\frac{(x-\mu)^2}{2\sigma^2}.
\]
(We used $\log(1/\sqrt{y})=-\tfrac12\log y$.)
\item Change the sign:
$-\log p(x)=\tfrac12\log(2\pi\sigma^2)+\dfrac{(x-\mu)^2}{2\sigma^2}$.
\item Take the expectation for $X\sim\N(\mu,\sigma^2)$. The first term is a
constant and stays as it is. In the second, the only random part is
$(X-\mu)^2$, whose expectation is \emph{by definition} the variance:
$\E[(X-\mu)^2]=\sigma^2$. Hence
\[
h=-\E[\log p(X)]=\tfrac12\log(2\pi\sigma^2)+\frac{\sigma^2}{2\sigma^2}
=\tfrac12\log(2\pi\sigma^2)+\tfrac12 .
\]
\item Put the $\tfrac12$ inside the logarithm by writing
$\tfrac12=\tfrac12\log e$ (since $\log e=1$), then
$\tfrac12\log A+\tfrac12\log e=\tfrac12\log(Ae)$:
\[
\boxed{\,h=\tfrac12\log(2\pi e\,\sigma^2)=\tfrac12\log(2\pi e)+\log\sigma\,}
\]
\end{enumerate}
Reading: the mean $\mu$ has disappeared (shifting a distribution does not
change its uncertainty), and only the volatility $\sigma$ matters, through
$\log\sigma$. Doubling volatility adds $\log2\approx0.69$ nat. Numerically,
for $\sigma=1$, $h=\tfrac12\log(2\pi e)=\tfrac12\log(17.08)\approx1.419$ nats.
\end{exemple}

Two important differences from the discrete case show up at once. First, $h$
can be negative: the uniform on $[0,1/2]$ has $h=-\log2<0$. Second, and this
is what matters for the paper, \textbf{$h$ depends on units}.

\begin{proposition}[Change of scale]
\label{prop:echelle}
For $a\neq0$, $h(aX+b)=h(X)+\log|a|$. In dimension $N$, for invertible $A$,
$h(AX+b)=h(X)+\log|\det A|$.
\end{proposition}

\begin{proof}
If $Y=AX+b$, the change-of-variables formula gives
$p_Y(y)=p_X(x)/|\det A|$ with $x=A^{-1}(y-b)$. Hence
$h(Y)=-\E[\log p_Y(Y)]=-\E[\log p_X(X)]+\log|\det A|=h(X)+\log|\det A|$.
\end{proof}

\begin{piege}
Expressing returns in per cent rather than decimals multiplies each coordinate
by $100$, and so adds $N\log100$ to $h$. The \emph{level} of $h$ therefore has
no absolute meaning; only its changes, and above all the \emph{divergences} of
the next section, do. This is why the paper builds all its indicators as
divergences.
\end{piege}

\begin{exercice}
Compute $h$ for $X\sim\N(0,\sigma^2)$ with $\sigma=1\,\%$, writing first
$\sigma=0.01$ and then $\sigma=1$ (returns in per cent). Check that the gap is
$\log100\approx4.61$.
\end{exercice}
\noindent\emph{Solution.} $h=\tfrac12\log(2\pi e)+\log\sigma$, that is
$1.419-4.605=-3.186$, then $1.419$. The gap is $\log 100$.

%======================================================================
\section{The Kullback--Leibler divergence}
\label{sec:kl}

\begin{definition}
For two densities $p$ and $q$ (with $q>0$ wherever $p>0$),
\[
\KL{p}{q}=\int p(x)\log\frac{p(x)}{q(x)}\,dx=\E_p\!\left[\log\frac{p(X)}{q(X)}\right].
\]
\end{definition}

\begin{intuition}
It is the extra surprise, on average, when you believe the world follows $q$
while it actually follows $p$. Another reading: when testing $p$ against $q$,
the log-likelihood ratio rises on average by $\KL{p}{q}$ per observation. A
divergence of one nat means that each observation multiplies the evidence in
favour of $p$ by $e$ on average.
\end{intuition}

\subsection{Three properties}

\begin{proposition}[Gibbs' inequality]
\label{prop:gibbs}
$\KL{p}{q}\ge0$, with equality if and only if $p=q$ almost everywhere.
\end{proposition}

\paragraph{The tool: Jensen's inequality.} A function $\varphi$ is
\emph{concave} if its graph always lies above its chords: for two points $a,b$
and a weight $t\in[0,1]$, $\varphi(ta+(1-t)b)\ge t\varphi(a)+(1-t)\varphi(b)$.
The logarithm is concave (its second derivative $-1/x^2$ is negative). Jensen
extends this chord inequality to any average: for any random variable $Z>0$,
\[
\E[\log Z]\ \le\ \log\E[Z]
\qquad\text{(``the mean of the logs is below the log of the mean'')}.
\]
Example: $Z$ equals $1$ or $4$ with probability $\tfrac12$. Then
$\E[\log Z]=\tfrac12(0+\log4)=\log2\approx0.69$, whereas
$\log\E[Z]=\log2.5\approx0.92$. Since the logarithm is \emph{strictly}
concave, equality holds only if $Z$ is constant: if $Z$ always equals $c$,
both sides equal $\log c$.

\begin{proof}
We proceed in four steps.

\emph{Step 1: flip the fraction.} Since $\log(a/b)=-\log(b/a)$,
\[
-\KL{p}{q}=-\E_p\!\left[\log\frac{p(X)}{q(X)}\right]
=\E_p\!\left[\log\frac{q(X)}{p(X)}\right].
\]
It is enough to show that this last number is $\le0$.

\emph{Step 2: apply Jensen.} Let $Z=q(X)/p(X)$, where $X$ follows $p$. Jensen
gives
\[
\E_p\!\left[\log\frac{q(X)}{p(X)}\right]\ \le\ \log\E_p\!\left[\frac{q(X)}{p(X)}\right].
\]

\emph{Step 3: compute the expectation on the right.} By definition of the
expectation under $p$, we integrate against the density $p$, over the region
where $p>0$ (elsewhere $X$ never falls):
\[
\E_p\!\left[\frac{q(X)}{p(X)}\right]
=\int_{p>0}p(x)\,\frac{q(x)}{p(x)}\,dx
=\int_{p>0}q(x)\,dx .
\]
The $p$'s cancel: what remains is the mass $q$ puts on the region where
$p>0$. This is a probability, so it is $\le1$ (exactly $1$ if $q$ puts no mass
outside this region).

\emph{Step 4: conclude.} The logarithm is increasing, so
$\log\int_{p>0}q\le\log1=0$. Chaining the steps,
\[
-\KL{p}{q}\ \le\ \log\int_{p>0}q\ \le\ 0,
\qquad\text{that is}\qquad \KL{p}{q}\ge0 .
\]

\emph{Equality case.} If $\KL{p}{q}=0$, both inequalities are equalities.
Equality in Jensen requires $Z=q/p$ to be constant, say $q(x)=c\,p(x)$
wherever $p>0$. Equality in step 4 requires $\int_{p>0}q=1$, that is
$c\int p=c=1$. Hence $q=p$. (``Almost everywhere'' only means that sets of
probability zero, such as a few isolated points, can be ignored, since they
change no integral.)
\end{proof}

\begin{exemple}[Check on a simple case]
A biased coin $p=(0.9,\,0.1)$ compared with a fair coin $q=(0.5,\,0.5)$:
\[
\KL{p}{q}=0.9\log\frac{0.9}{0.5}+0.1\log\frac{0.1}{0.5}
=0.9\times0.588+0.1\times(-1.609)=0.529-0.161=0.368 .
\]
One of the two terms is negative, but the weighted sum is positive: exactly
what the proposition guarantees. The negative term is always ``dominated''
because it carries the small weight $0.1$.
\end{exemple}

\begin{remarque}[A proof without Jensen]
One can also use the elementary inequality $\log x\le x-1$ (the graph of the
logarithm lies below its tangent at $x=1$, with equality only at $x=1$). With
$x=q/p$: $\E_p[\log(q/p)]\le\E_p[q/p-1]=\int_{p>0}q-1\le0$. Same idea,
written with a tangent instead of a chord.
\end{remarque}

\begin{proposition}[Invariance under change of variables]
\label{prop:invariance}
If $Y=\phi(X)$ with $\phi$ bijective and differentiable, and $p_Y,q_Y$ are the
laws of $Y$ when $X\sim p$ or $X\sim q$, then $\KL{p_Y}{q_Y}=\KL{p}{q}$.
\end{proposition}

\begin{proof}
$p_Y(y)=p(x)/|J_\phi(x)|$ and $q_Y(y)=q(x)/|J_\phi(x)|$ with the \emph{same}
Jacobian. In the ratio $p_Y/q_Y$, the Jacobian cancels.
\end{proof}

\begin{papier}
This is the property that lets equity returns, rate changes and exchange rates
sit in the same matrix (entropic CISS), and nats be compared across markets.
\end{papier}

\paragraph{Not a distance.} In general $\KL{p}{q}\neq\KL{q}{p}$, and the
triangle inequality fails: hence the word ``divergence''. For example, for two
centred Gaussians with variances $\sigma_1^2$ and $\sigma_0^2$, we show
(section~\ref{sec:scenarios}) that $\KL{p_1}{p_0}=\tfrac12(v^2-1-2\log v)$ with
$v=\sigma_1/\sigma_0$. For $v=2$ this gives $0.807$, and for $v=1/2$, $0.318$.

\subsection{The most useful case: two Gaussians with the same covariance}

\begin{proposition}
\label{prop:klmoy}
If $p_1=\N(\mu_1,\Sigma)$ and $p_0=\N(\mu_0,\Sigma)$, then
$\KL{p_1}{p_0}=\tfrac12\,\Delta^\top\Sigma^{-1}\Delta$ with
$\Delta=\mu_1-\mu_0$.
\end{proposition}

\begin{proof}
The normalising constants are identical and cancel. What remains is
\[
\log\frac{p_1(x)}{p_0(x)}=-\tfrac12(x-\mu_1)^\top\Sigma^{-1}(x-\mu_1)
+\tfrac12(x-\mu_0)^\top\Sigma^{-1}(x-\mu_0).
\]
Under $p_1$, write $x-\mu_0=(x-\mu_1)+\Delta$ and expand: the cross term
$\Delta^\top\Sigma^{-1}(X-\mu_1)$ has zero expectation, and the two quadratic
forms in $X-\mu_1$ cancel. What remains is $\tfrac12\Delta^\top\Sigma^{-1}\Delta$.
\end{proof}

The quantity $\sqrt{\Delta^\top\Sigma^{-1}\Delta}$ is the \emph{Mahalanobis
distance}: the shift measured in standard deviations, taking correlations into
account. In one dimension, $\Delta=k\sigma$ gives
\[
\boxed{\;k\text{ standard deviations cost }k^2/2\text{ nats, and }\eta\text{ nats
buy }\sqrt{2\eta}\text{ standard deviations.}\;}
\]
A $3\sigma$ shock costs $4.5$ nats; one nat buys $\sqrt2\approx1.41$ standard
deviations. This conversion rule runs through the whole paper.

%======================================================================
\section{Why the Gaussian is the zero point}
\label{sec:maxent}

\begin{proposition}[Maximum entropy]
\label{prop:maxent}
Among all densities on $\R^N$ with mean $\mu$ and covariance $\Sigma\succ0$,
the Gaussian $q=\N(\mu,\Sigma)$ is the unique maximum-entropy density.
\end{proposition}

\paragraph{What needs to be shown.} Fix a mean and a covariance. Many
distributions have them: a Gaussian, a suitably chosen uniform, a Laplace, a
fat-tailed law\ldots{} The proposition says that among all of them, the
Gaussian has the highest entropy. To show it, we compare an arbitrary $p$ with
the Gaussian $q$ with the same moments, using the divergence $\KL{p}{q}$,
which we already know is non-negative (Gibbs).

\paragraph{Ingredient 1: the log of a Gaussian is a degree-two polynomial.}
Start in one dimension. The density of $q=\N(\mu,\sigma^2)$ is
$q(x)=\frac{1}{\sqrt{2\pi\sigma^2}}\exp\!\bigl(-\frac{(x-\mu)^2}{2\sigma^2}\bigr)$.
The logarithm ``undoes'' the exponential (example~\ref{ex:unif-norm}):
\[
\log q(x)=\underbrace{-\tfrac12\log(2\pi\sigma^2)}_{\text{constant}}
\;-\;\frac{(x-\mu)^2}{2\sigma^2}.
\]
Expanding $(x-\mu)^2=x^2-2\mu x+\mu^2$ shows it has the form
\[
\log q(x)=a+b\,x+c\,x^2,
\qquad a=-\tfrac12\log(2\pi\sigma^2)-\frac{\mu^2}{2\sigma^2},\quad
b=\frac{\mu}{\sigma^2},\quad c=-\frac{1}{2\sigma^2}.
\]
This is what ``$\log q$ is quadratic'' means: a degree-two polynomial in $x$,
with no $x^3$, $x^4$, etc. This property is \emph{specific} to the Gaussian.
For a Laplace law, $\log p(x)=\text{const}-|x|/b$ is not a polynomial; for a
Student law, $\log p$ contains $\log(1+x^2/\nu)$. In dimension $N$ it is the
same:
\[
\log q(x)=-\tfrac12\bigl[N\log2\pi+\log\det\Sigma\bigr]
-\tfrac12(x-\mu)^\top\Sigma^{-1}(x-\mu),
\]
a constant plus a quadratic form in $x$.

\paragraph{Ingredient 2: the expectation of a degree-two polynomial depends
only on the first two moments.} By linearity of expectation, for any variable
$X$ (with any law),
\[
\E[a+bX+cX^2]=a+b\,\E[X]+c\,\E[X^2].
\]
But $\E[X]=\mu$ and $\E[X^2]=\sigma^2+\mu^2$ are fixed by the mean and
variance. \textbf{All laws with mean $\mu$ and variance $\sigma^2$ therefore
give the same value to $\E[\log q(X)]$}, Gaussian or not. This would not hold
for a function containing $x^4$: its expectation would involve the kurtosis,
which differs from one law to another.

\paragraph{Consequence.} If $p$ has the same mean and variance as $q$, $p$ can
be replaced by $q$ in the expectation of $\log q$:
\[
\E_p[\log q(X)]=\E_q[\log q(X)]=-h(q).
\]
The last equality is the definition of the entropy of $q$.

\begin{proof}
Put the pieces together. By definition of the divergence, then of entropy,
\[
\KL{p}{q}=\E_p[\log p(X)]-\E_p[\log q(X)]
=-h(p)-\E_p[\log q(X)].
\]
Replace $\E_p[\log q(X)]$ by $-h(q)$ (consequence above):
\[
\KL{p}{q}=-h(p)+h(q)=h(q)-h(p).
\]
By Gibbs, $\KL{p}{q}\ge0$, so $h(p)\le h(q)$. Equality holds if and only if
$\KL{p}{q}=0$, that is $p=q$: the Gaussian is the unique maximum.
\end{proof}

\begin{exemple}[Numerical check, mean $0$ and variance $1$]
\leavevmode
\begin{itemize}
\item \emph{Gaussian} $\N(0,1)$: $h=\tfrac12\log(2\pi e)\approx1.419$.
\item \emph{Uniform on $[-\sqrt3,\sqrt3]$} (its variance is
$(2\sqrt3)^2/12=1$): $h=\log(2\sqrt3)\approx1.242$.
\item \emph{Laplace} with density $\frac{1}{2b}e^{-|x|/b}$ and $b=1/\sqrt2$
(variance $2b^2=1$): $h=1+\log(2b)\approx1.347$.
\end{itemize}
The Gaussian does have the highest entropy. Let us also check ingredient~2 on
the uniform: $\log q(x)=-\tfrac12\log2\pi-\tfrac12x^2$, so
$\E_p[\log q]=-\tfrac12\log2\pi-\tfrac12\E_p[X^2]
=-\tfrac12\log2\pi-\tfrac12=-1.419=-h(q)$, although $p$ is not Gaussian.
Finally $\KL{p}{q}=1.419-1.242=0.177$ nat: the ``entropy deficit'' of the
uniform, what the Gaussian does not see.
\end{exemple}

\begin{intuition}
Fixing a mean and a variance, and assuming nothing else, leads to the
Gaussian. Any other law with the same moments ``knows more'': it contains
additional information (skewness, fat tails, nonlinear dependence). Measuring
the departure from the Gaussian therefore measures this information.
\end{intuition}

\begin{remarque}
One can also find the Gaussian with a Lagrangian: maximising $-\int p\log p$
subject to $\int p=1$, $\int xp=\mu$, $\int xx^\top p=\Sigma+\mu\mu^\top$ gives
a first-order condition in which $\log p$ is quadratic, hence $p$ Gaussian.
This route explains the form of the result but proves neither that it is a
maximum nor that it is unique; the divergence proof does both in three lines.
\end{remarque}

\paragraph{Where does the variance come from? (the ``bridge'' of the online appendix).}
Maximum entropy takes the variance as given without saying where it comes
from. The paper proposes measuring time in traded volume. Over an interval,
let $V$ be volume and $Q$ net order flow. Two assumptions: (i) the change in
log price is $\lambda Q$ (linear Kyle-type impact); (ii) given $V$, all we
know about $Q$ is $\E[Q\mid V]=m$ and $\operatorname{Var}(Q\mid V)=\varsigma^2V$.
Proposition~\ref{prop:maxent} in one dimension then gives
\[
\Delta X\mid V\sim\N(\lambda m,\ \lambda^2\varsigma^2V).
\]
Unconditionally, $\Delta X$ is a \emph{mixture} of Gaussians with random
variances, hence fat-tailed. The moral: fat tails in calendar time do not
contradict maximum entropy; they can come from a random clock. No empirical
result depends on this bridge: it only justifies the Gaussian as the
reference.

\begin{exercice}[A mixture has fat tails]
Let $X=\sqrt{V}\,Z$ with $Z\sim\N(0,1)$ independent of $V>0$. Show that the
kurtosis of $X$ is $3\,\E[V^2]/\E[V]^2\ge3$, with equality if and only if $V$
is constant.
\end{exercice}
\noindent\emph{Solution.} $\E[X^2]=\E[V]$ and $\E[X^4]=\E[V^2]\E[Z^4]=3\E[V^2]$.
The ratio is $\E[X^4]/\E[X^2]^2=3\E[V^2]/\E[V]^2$, and $\E[V^2]\ge\E[V]^2$ by
Jensen (or because the variance of $V$ is non-negative).

%======================================================================
\section{The entropy of a multivariate Gaussian}
\label{sec:gauss}

\begin{proposition}
\label{prop:hgauss}
For $q=\N(\mu,\Sigma)$ on $\R^N$,
\[
h(q)=\tfrac12\log\bigl((2\pi e)^N\det\Sigma\bigr).
\]
\end{proposition}

\begin{proof}
Start from the expression of $\log q$ above:
$h(q)=\tfrac N2\log2\pi+\tfrac12\log\det\Sigma
+\tfrac12\E[(X-\mu)^\top\Sigma^{-1}(X-\mu)]$.
For the last expectation, use \textbf{the trace trick}: a scalar equals its
trace, and the trace is invariant under cyclic permutation
($\tr(ABC)=\tr(CAB)$). Hence
\[
\E\bigl[(X-\mu)^\top\Sigma^{-1}(X-\mu)\bigr]
=\E\bigl[\tr\bigl(\Sigma^{-1}(X-\mu)(X-\mu)^\top\bigr)\bigr]
=\tr\bigl(\Sigma^{-1}\Sigma\bigr)=\tr(I_N)=N .
\]
So $h(q)=\tfrac12\bigl[N\log2\pi+\log\det\Sigma+N\bigr]$.
\end{proof}

\begin{remarque}
The trace trick comes back several times. Remember: for any $X$ with
covariance $\Sigma$ and any matrix $A$,
$\E[(X-\mu)^\top A(X-\mu)]=\tr(A\Sigma)$, \emph{whatever the law of $X$}.
Only the first two moments enter.
\end{remarque}

\subsection{Separating volatilities and correlations}

With $\Sigma=DRD$, multiplicativity of the determinant gives
$\det\Sigma=(\det D)^2\det R=\bigl(\prod_i\sigma_i\bigr)^2\det R$, so
\begin{equation}
\label{eq:plafond}
h(q)=\underbrace{\tfrac N2\log(2\pi e)}_{\text{constant}}
+\underbrace{\sum_i\log\sigma_i}_{\text{scale}}
+\underbrace{\tfrac12\log\det R}_{\text{dependence}} .
\end{equation}

\begin{proposition}[Correlation destroys entropy]
\label{prop:hadamard}
For every correlation matrix $R$, $\det R\le1$, with equality if and only if
$R=I$. The dependence term is therefore always negative or zero.
\end{proposition}

\begin{proof}
$R$ is symmetric positive semi-definite, with eigenvalues
$\lambda_1,\dots,\lambda_N\ge0$ and $\sum_k\lambda_k=\tr R=N$ (the diagonal of
$R$ contains only $1$s). By the arithmetic--geometric mean inequality,
\[
\det R=\prod_k\lambda_k\le\Bigl(\tfrac1N\sum_k\lambda_k\Bigr)^N=1,
\]
with equality if and only if every $\lambda_k$ equals $1$, that is $R=I$.
(This is a special case of Hadamard's inequality.)
\end{proof}

\begin{papier}
Equation~\eqref{eq:plafond} is the paper's ``ceiling'': what we would get if
returns were Gaussian. It also explains why the scale channel is a nuisance:
in a crisis, volatilities rise, so $h(q)$ \emph{rises} (by $5.97$ nats on the
main panel), whereas we would like a crisis to show up as a
\emph{constraint}. Hence the idea of removing scale.
\end{papier}

%======================================================================
\section{Mutual information and coupling}
\label{sec:couplage}

\begin{definition}
The mutual information of two vectors $X$ and $Y$ is
$I(X;Y)=\KL{p_{XY}}{p_X\otimes p_Y}$: the gap between the joint law and the
law we would have if $X$ and $Y$ were independent. For $N$ variables, the
\emph{total correlation} (or multi-information, Watanabe 1960) is
$\KL{p}{\prod_ip_i}$.
\end{definition}

By Gibbs, $I\ge0$, with equality if and only if the variables are
independent. By proposition~\ref{prop:invariance}, $I$ does not depend on
units.

\begin{proposition}[The Gaussian case]
\label{prop:migauss}
If $X\sim\N(0,\Sigma)$ with $\Sigma=DRD$, the total correlation is
\[
C(R)=\KL{\N(0,\Sigma)}{\N(0,D^2)}=-\tfrac12\log\det R .
\]
For two variables with correlation $\rho$: $I=-\tfrac12\log(1-\rho^2)$.
\end{proposition}

\begin{proof}
The marginals of a Gaussian are Gaussian, so $\prod_ip_i=\N(0,D^2)$. For a
Gaussian vector, $\KL{p}{\prod_ip_i}=\sum_ih(p_i)-h(p)$ (same argument as
proposition~\ref{prop:maxent}: $\log\prod_ip_i$ is quadratic and $p$ has the
same variances). With~\eqref{eq:plafond}:
$\sum_ih(p_i)=\tfrac N2\log(2\pi e)+\sum_i\log\sigma_i$, and the difference
with $h(p)$ is $-\tfrac12\log\det R$. For $N=2$, $\det R=1-\rho^2$.
\end{proof}

\begin{exemple}
For $\rho=0.9$, $I=-\tfrac12\log0.19=0.83$ nat. For $\rho=0.3$,
$I=0.047$ nat: mutual information rises very fast as $\rho$ approaches $1$.
\end{exemple}

\paragraph{Constant correlation.} If every correlation equals $\rho$,
$R=(1-\rho)I+\rho\,\mathbf 1\mathbf 1^\top$. The vector $\mathbf 1$ is an
eigenvector with eigenvalue $1+(N-1)\rho$ (the ``market mode''), and every
vector orthogonal to $\mathbf 1$ has eigenvalue $1-\rho$. Hence
\[
\det R=(1-\rho)^{N-1}\bigl(1+(N-1)\rho\bigr),\qquad
C=-\tfrac12\bigl[(N-1)\log(1-\rho)+\log(1+(N-1)\rho)\bigr].
\]
For $N=20$: $C=1.97$ nats if $\rho=0.26$ and $C=2.69$ nats if $\rho=0.32$.
These are the average calm and stress correlations of the main panel; the
gap, $0.72$ nat, gives the order of magnitude of the $0.99$-nat rise measured
on the actual matrix.

\paragraph{Effective number of modes.} A related spectral reading: with
$\tilde\lambda_k=\lambda_k/N$ (which sum to $1$, like probabilities), define
$N^{\text{eff}}=\exp\bigl(-\sum_k\tilde\lambda_k\log\tilde\lambda_k\bigr)$.
It is the exponential of the entropy of the spectrum: $N^{\text{eff}}=N$ if
$R=I$, and $N^{\text{eff}}\to1$ if a single mode carries all the variance.
With the constant correlation above, $N^{\text{eff}}$ goes from $14.6$ to
$12.8$ (on the actual matrix, the paper finds $13.1$ then $11.4$).

\begin{intuition}
$C$ and $N^{\text{eff}}$ say the same thing in two ways: under stress, there
are fewer independent sources of risk. But $C=-\tfrac12\sum_k\log\lambda_k$
\emph{weighs the small eigenvalues heavily} (the $\log$ blows up as
$\lambda_k\to0$). And these are the ones that collapse when assets couple.
This is the explanation offered for the robustness of $C$ compared with
measures that look only at the largest eigenvalue (absorption ratio).
\end{intuition}

\begin{exercice}
Show that for small $\rho$, $C\approx\tfrac14N(N-1)\rho^2$.
\end{exercice}
\noindent\emph{Solution.} $\log(1-\rho)\approx-\rho-\rho^2/2$ and
$\log(1+(N-1)\rho)\approx(N-1)\rho-(N-1)^2\rho^2/2$. The first-order terms
cancel; at second order,
$C\approx\tfrac12\bigl[(N-1)\rho^2/2+(N-1)^2\rho^2/2\bigr]=\tfrac14N(N-1)\rho^2$.
Consequence: $C$ grows like $N^2$; comparing panels of different sizes
requires normalisation.

%======================================================================
\section{The three channels and the structural index}
\label{sec:canaux}

\subsection{The entropy deficit, or shape channel}

\begin{definition}
Let $q$ be the Gaussian with the same mean and covariance as the true law $p$.
The \emph{entropy deficit} is $\DD=\KL{p}{q}$.
\end{definition}

\begin{proposition}
\label{prop:deficit}
$\DD=h(q)-h(p)\ge0$, with equality if and only if $p$ is Gaussian.
\end{proposition}

\begin{proof}
This is exactly the proof of proposition~\ref{prop:maxent} read backwards:
$\E_p[\log q]=-h(q)$ because $\log q$ is quadratic and $p$, $q$ share the same
moments, so $\KL{p}{q}=-h(p)+h(q)$.
\end{proof}

\begin{piege}
The identity $\KL{p}{q}=h(q)-h(p)$ is false in general. It holds only because
$q$ is Gaussian \emph{and} has the same first two moments as $p$.
\end{piege}

\subsection{The exact decomposition}

Combining~\eqref{eq:plafond} and proposition~\ref{prop:deficit} gives the
central equation of the framework (eq.~(1) of the paper):
\begin{equation}
\label{eq:trois}
h(p)=\tfrac N2\log(2\pi e)
+\underbrace{\sum_i\log\sigma_i}_{\text{scale}}
+\underbrace{\tfrac12\log\det R}_{\text{dependence }=-C}
-\underbrace{\DD}_{\text{shape}} .
\end{equation}
There is no approximation: it is an identity.

\subsection{The structural index is a divergence}

We remove the scale channel and keep the other two: $\JJ=\DD+C$. This choice
might look like a hack. It is not:

\begin{proposition}
\label{prop:J}
Let $q_i=\N(\mu_i,\sigma_i^2)$ be the Gaussians with the same means and
volatilities as each asset. Then
\[
\JJ=\DD+C=\KL{p}{\textstyle\prod_iq_i}\ \ge0 .
\]
\end{proposition}

\begin{proof}
$\log\prod_iq_i(x)=-\sum_i\bigl[\tfrac12\log(2\pi\sigma_i^2)
+\tfrac{(x_i-\mu_i)^2}{2\sigma_i^2}\bigr]$. Under $p$,
$\E[(X_i-\mu_i)^2]=\sigma_i^2$, so
\[
\KL{p}{\textstyle\prod_iq_i}=-h(p)+\tfrac N2\log(2\pi e)+\sum_i\log\sigma_i .
\]
Substitute~\eqref{eq:trois} for $h(p)$: the constant and the scale term
cancel, leaving $-\tfrac12\log\det R+\DD=C+\DD$.
\end{proof}

Being a divergence gives $\JJ$ three properties for free: it is
non-negative; it is zero if and only if returns are \emph{Gaussian and
independent} (a well-defined reference); and it depends neither on units nor
on the level of volatility.

\paragraph{The chain rule.} We also have the general identity
\[
\KL{p}{\textstyle\prod_iq_i}=\sum_i\KL{p_i}{q_i}+\KL{p}{\textstyle\prod_ip_i},
\]
(marginal deficits + full multi-information). The paper estimates the first
term asset by asset and replaces the second by its Gaussian part $C$.
Nonlinear tail dependence is therefore left out: the estimated $\JJ$
\emph{understates} the true divergence (a limit stated in the paper).

\subsection{How the marginal deficit is estimated}

For a standardised variable with skewness $S$ and excess kurtosis $K$, close
to Gaussian, the Edgeworth expansion gives
\[
\DD\approx\frac{S^2}{12}+\frac{K^2}{48}.
\]
This helps understanding (a skewness channel and a tail channel, additive)
but is unusable as an estimator on daily returns: the estimated $K$ is
dominated by a few extreme days, and squaring it amplifies the problem (the
paper reports a factor-69 overstatement for a Student~$t$ with 4 degrees of
freedom). Instead we use the Hyvärinen (1998) approximation, of the form
\[
\DD\approx k_1\bigl(\E[G_1(Z)]\bigr)^2+k_2\bigl(\E[G_2(Z)]-\E[G_2(\nu)]\bigr)^2,
\qquad\nu\sim\N(0,1),
\]
where $G_1$ is a bounded odd function (it captures skewness) and $G_2$ a
slowly growing even function (it captures tails). Since $G_1,G_2$ grow more
slowly than any polynomial, the estimator is much less sensitive to extreme
values.

\paragraph{An important property: $\JJ$ is even in skewness.} The reference
$\N(0,1)$ is symmetric, so replacing $Z$ by $-Z$ does not change
$\KL{p_Z}{\varphi}$. Edgeworth shows it too: $S$ enters squared. \textbf{A
crash and its mirror image (a rally) have the same $\JJ$.} The index measures
the \emph{size} of information, never its direction. Hence the need for an
outside postulate on the sign (next section).

%======================================================================
\section{Dynamics: a theorem and two postulates}
\label{sec:dynamique}

\subsection{The theorem: information lowers uncertainty on average}

Let $h(X\mid Y=y)$ be the entropy of the conditional law of $X$ given $Y=y$,
and $\E_Y[h(X\mid Y)]$ its average over the values of the signal.

\begin{proposition}
\label{prop:cond}
$\E_Y[h(X\mid Y)]=h(X)-I(X;Y)\le h(X)$.
\end{proposition}

\begin{proof}
By definition, $p(x\mid y)=p(x,y)/p(y)$, so
$\E_Y[h(X\mid Y)]=-\E[\log p(X,Y)]+\E[\log p(Y)]=h(X,Y)-h(Y)$. Moreover
$I(X;Y)=\E\bigl[\log\tfrac{p(X,Y)}{p(X)p(Y)}\bigr]=h(X)+h(Y)-h(X,Y)$. Adding,
$\E_Y[h(X\mid Y)]=h(X)-I(X;Y)$, and $I\ge0$ by Gibbs.
\end{proof}

\begin{piege}
The theorem holds \emph{on average}. A particular realisation of the signal
can increase entropy. Example: $X$ equals $a$ with probability $0.99$, so
$H(X)=0.056$ nat. With probability $0.98$ the signal removes all doubt
($H=0$); with probability $0.02$ it makes the two outcomes equally likely
($H=\log2=0.693$, twelve times more than before). The average,
$0.02\times0.693=0.014$, is indeed below $0.056$. This is why the framework
predicts jumps that are \emph{mainly}, not exclusively, one-sided.
\end{piege}

\begin{exemple}[A noisy Gaussian signal]
Let $X\sim\N(0,\sigma^2)$ be tomorrow's return and $Y=X+\varepsilon$ with
$\varepsilon\sim\N(0,\sigma_\varepsilon^2)$ independent. The pair $(X,Y)$ is
Gaussian with $\rho^2=\sigma^2/(\sigma^2+\sigma_\varepsilon^2)$, so by
proposition~\ref{prop:migauss}
\[
I(X;Y)=-\tfrac12\log(1-\rho^2)=\tfrac12\log\Bigl(1+\frac{\sigma^2}{\sigma_\varepsilon^2}\Bigr).
\]
With $\sigma=1\,\%$: a sharp signal ($\sigma_\varepsilon=0.5\,\%$) brings
$0.80$ nat, an average signal ($1\,\%$) $0.35$ nat, a rumour ($3\,\%$)
$0.05$ nat. A factor of fifteen separates clear news from a rumour.
\end{exemple}

\subsection{Postulate 1: relaxation (Ornstein--Uhlenbeck process)}

Without news, we postulate that $\JJ$ reverts to its normal level $\bar\JJ$:
\[
d\JJ_t=-\kappa(\JJ_t-\bar\JJ)\,dt+\sigma_\JJ\,dW_t .
\]

\paragraph{The half-life.} Taking expectations, the Brownian term vanishes and
$m(t)=\E[\JJ_t]-\bar\JJ$ satisfies $m'=-\kappa m$, so $m(t)=m(0)e^{-\kappa t}$.
The gap halves when $e^{-\kappa t}=1/2$, that is
\[
t_{1/2}=\frac{\ln2}{\kappa}.
\]
In discrete time (daily data), the model becomes an AR(1),
$\JJ_{t+1}-\bar\JJ=\phi(\JJ_t-\bar\JJ)+\varepsilon_t$ with $\phi=e^{-\kappa}$,
and the half-life is $\ln2/(-\ln\phi)$ days. This is how the $121$-day
half-life (against $92$ under the null) is estimated.

\paragraph{Why this process?} Doob's theorem (1942) says that the
Ornstein--Uhlenbeck process is the only continuous process that is Markovian,
Gaussian and stationary. If we want a process with no memory beyond the
current state, a long-run level, and noise about which we assume nothing
beyond its variance (maximum entropy again), there is no choice. We then add
the jumps caused by news:
\[
d\JJ_t=-\kappa(\JJ_t-\bar\JJ)\,dt+\sigma_\JJ\,dW_t+dZ_t ,
\]
where $Z_t=\sum_{j:\,t_j\le t}z_j$ is the sum of positive jumps up to $t$
(Barndorff-Nielsen and Shephard, 2001, class). What remains arbitrary:
constant $\kappa$, linear mean reversion, and the fact that the Brownian term
could make $\JJ$ negative.

\subsection{Postulate 2: the sign}

Since $\JJ$ does not see direction (section~\ref{sec:canaux}), we add a
minimal assumption: jumps come with negative skewness, justified by the
leverage effect (falls raise volatility more than rises of the same size).
This is hypothesis H5.

\subsection{Why dependence and tails should rise together}

Major news hits every asset on the same day. On that day, they move together
(coupling $C$ rises) and each records an unusual loss (its deficit $\DD$
rises). The two channels describe the same day seen from two places: the
cross-section and each asset's time series. This is hypothesis H6, which we
will see the null reproduces.

%======================================================================
\section{Inference: why and how the placebo nulls work}
\label{sec:temoin}

This is the most important section for understanding the results of the
journal version. It explains why estimated entropy measures are biased and
move on their own, how a resampling null corrects for it, why a second,
volatility-preserving null is needed for claims about dependence, and how the
size and power of the test are checked.

\subsection{The bias of a convex function of estimated moments}

\begin{lemme}[Jensen applied to an estimator]
If $\hat\theta$ is an (approximately) unbiased estimator of $\theta$ and $f$
is convex, then $\E[f(\hat\theta)]\ge f(\theta)$.
\end{lemme}

\noindent $C=-\tfrac12\log\det R$ is a convex function of $R$ ($-\log\det$ is
convex on positive definite matrices), and $\DD$ is non-negative and zero only
at the Gaussian point. Both are therefore \textbf{biased upward}: estimation
noise looks like information.

\subsection{How large is the bias of coupling? (Proposition 3 of the paper)}

Take $N$ \emph{independent} Gaussian assets, so the true coupling is $C=0$,
and $T$ observations. The likelihood-ratio statistic for testing $H_0:R=I$
(``sphericity'' of the correlation matrix) is
\[
\Lambda=-T\log\det\hat R .
\]
Classical asymptotic theory says that $-2\log(\text{likelihood ratio})$ is
approximately $\chi^2$ with as many degrees of freedom as free parameters
restricted by $H_0$. Here $H_0$ fixes the $N(N-1)/2$ off-diagonal
correlations, so $\Lambda\approx\chi^2_{N(N-1)/2}$. Bartlett (1950) showed that
replacing $T$ by $T'=T-1-(2N+5)/6$ makes the approximation much better in
small samples. Since a $\chi^2_k$ variable has mean $k$,
\[
\E\bigl[-\log\det\hat R\bigr]\approx\frac{N(N-1)}{2T'}
\quad\Longrightarrow\quad
\E\bigl[C(\hat R)\bigr]\approx b(N,T)=\frac{N(N-1)}{4\bigl(T-1-(2N+5)/6\bigr)} .
\]

\begin{exemple}[The paper's panels]
For $N=20$, $T=252$: $T'=251-7.5=243.5$ and $b=380/974=\BIASTWENTY$~nat. For
$N=48$, $T=252$: $T'=251-16.8=234.2$ and $b=2256/936.8=\BIASFORTYEIGHT$~nats.
The simulation of the paper (table~2) gives \BIASTWENTYMC~nat for $N=20$: the
formula is accurate. The bias is as large as a third of the effect we want
to measure on stocks ($0.99$ nat).
\end{exemple}

\begin{intuition}
Why fourth moments barely matter: for two independent variables with finite
fourth moments, $\sqrt T\hat\rho\to\N(0,1)$ whatever their marginal laws,
because $\operatorname{Var}(XY)=\operatorname{Var}(X)\operatorname{Var}(Y)$
when $X$ and $Y$ are independent. The bias formula only uses the variance of
$\hat\rho$, so it holds for Student-$t_5$ data too (the simulation gives
\BIASTWENTYT).
\end{intuition}

\paragraph{Two practical consequences.} (1) Coupling levels on panels of
different size are not comparable: subtract $b(N,T)$ first. The
bias-corrected effective number of modes is $N\exp\{-2[C(\hat R)-b]/N\}$.
(2) A \emph{difference} between two windows of the same length (stress minus
calm) has the bias removed to first order. This is why every test in the
paper is on a gap, not on a level.

\begin{piege}
Shrinkage (Ledoit--Wolf) pulls $\hat R$ towards $I$ and therefore removes most
of the bias (\BIASTWENTYLW~nat in the simulation). It does not remove the
\emph{dynamic} artefact of the next subsection.
\end{piege}

\subsection{The mechanical sawtooth}

When an extreme day enters the window, the estimated moments jump, and a
convex function of these moments jumps \emph{upward whatever the sign of the
day}. As the day ages (exponential weighting) or leaves the window, the index
decays slowly.

\paragraph{A short calculation.} Take a window of $T$ standardised returns
whose sample correlation is exactly $I$ (independent assets), and add one day
on which every asset moves by the same amount $a$ standard deviations. Ignoring
the change in the means, the new second-moment matrix is $I+x\mathbf 1\mathbf
1^\top$ with $x=a^2/T$, and the new correlation matrix is
$R'=(I+x\mathbf 1\mathbf 1^\top)/(1+x)$. By the matrix determinant lemma,
$\det(I+x\mathbf 1\mathbf 1^\top)=1+Nx$, so
\[
\Delta C=-\tfrac12\log\det R'=\tfrac12\bigl[N\log(1+x)-\log(1+Nx)\bigr]
\approx\frac{N(N-1)}{4}\,x^2\quad(x\text{ small}).
\]
The first-order terms cancel and what remains is positive: a single common day
raises coupling, whatever its sign. For $N=20$, $T=252$ and a 10-standard-deviation
day ($x=0.40$), $\Delta C=\tfrac12[20\log1.40-\log8.94]=2.2$~nats. The day then
leaves the estimate only at the rate at which the window forgets: with
exponential weights of half-life 60 days, its weight halves every 60 days.

On i.i.d.\ Gaussian data the paper's pipeline finds $37$ positive jumps of
$\JJ$, no negative one, and an apparent half-life of $35$ days, with no event
at all.

\subsection{The resampling test and why its $p$-value is valid}

\paragraph{Construction.} Let $S(X)$ be a statistic computed by the
\emph{whole} pipeline on the $T\times N$ panel $X$ (rolling estimation, stress
regime, jump detection, gap). A null is a way of generating panels $X^*$ that
keep some features of $X$ and destroy others. We draw $B=200$ of them and
compute $S(X^*_1),\dots,S(X^*_B)$.

\paragraph{The $p$-value.} In the predicted direction,
\[
\hat p=\frac{1+\#\{b:\ S(X^*_b)\ge S(X)\}}{B+1}.
\]
If, under the null hypothesis, the observed panel is just one more draw from
the same mechanism, then $S(X),S(X^*_1),\dots,S(X^*_B)$ are
\emph{exchangeable}: the observed value is equally likely to have any rank
among the $B+1$ values. The probability that it ranks in the top $k$ is
$k/(B+1)$, so $P(\hat p\le\alpha)\le\alpha$ exactly, without any asymptotic
argument. The ``$+1$'' is what makes this exact; with $B=200$ the smallest
possible $p$-value is $1/201\approx0.005$.

\begin{piege}
Reporting $z=(S(X)-\bar S^*)/s_{S^*}$ (``six standard deviations above the
null'') is a useful effect size, but it is not a $p$-value: the null
distribution need not be normal. The paper reports $\hat p$ for inference and
$z$ only as a description.
\end{piege}

\subsection{Null 1: i.i.d.\ resampling of days}

Draw rows of $X$ with replacement. Each row (one day, $N$ returns) is intact,
so the null keeps each asset's fat tails, each day's cross-sectional
correlation and common extreme days. It destroys \emph{all} dynamics: there
is no volatility clustering, no persistence, no crisis that lasts. It answers
the question ``beyond the fact that turbulent periods contain big days, does
the measure detect any dynamics?''

\subsection{Null 2: constant-correlation GARCH (filtered historical simulation)}

\paragraph{Why a second null.} Under the i.i.d.\ null, a ``stress'' period is
a random cluster of big days. In real data, stress comes with volatility
clustering: high-volatility days follow each other for weeks. Could
volatility clustering alone, with correlations that never change, produce the
coupling gap? Two mechanisms make this plausible. Forbes and Rigobon (2002)
showed that measured correlation rises mechanically when the variance of a
common factor rises. And our stress regime is defined by high volatility, so
it selects precisely the windows where this happens.

\paragraph{Construction.} For each asset $i$, fit a GJR-GARCH(1,1):
\[
r_{i,t}=\mu_i+\varepsilon_{i,t},\qquad \varepsilon_{i,t}=\sqrt{h_{i,t}}\,z_{i,t},\qquad
h_{i,t}=\omega_i+\bigl(\alpha_i+\gamma_i\mathbf 1\{\varepsilon_{i,t-1}<0\}\bigr)\varepsilon_{i,t-1}^2+\beta_ih_{i,t-1}.
\]
The term $\gamma_i$ makes volatility react more to falls than to rises (the
leverage effect). Estimate by maximising the Gaussian likelihood (quasi-MLE,
consistent even if $z$ is not Gaussian). Compute the standardised residuals
$\hat z_{i,t}=\hat\varepsilon_{i,t}/\sqrt{\hat h_{i,t}}$ and store them as a
$T\times N$ matrix $\hat Z$.

To simulate: draw rows of $\hat Z$ with replacement (as in null 1), then
rebuild each asset's returns by running its own GARCH recursion forward with
the drawn shocks. The result has three properties:
\begin{enumerate}
\item each asset keeps its own volatility clustering and leverage effect;
\item all assets receive the shocks of the same historical day, so a common
large day still hits everyone, and volatility bursts are common;
\item the correlation of the shocks $z_t$ is the same every day: the
\emph{conditional} correlation is constant (Bollerslev's CCC model).
\end{enumerate}
So volatility moves as in a GARCH world and correlations do not. A coupling
gap beyond this null is a genuine dependence dynamic.

\begin{intuition}
Null 1 says: ``it is not just big days''. Null 2 says: ``it is not just
volatility''. The 21-day block null that the earlier version rejected kept
\emph{both} volatility and correlation dynamics over a month, i.e.\ the crisis
itself. Null 2 keeps the volatility part only, which is exactly the
counterfactual the coupling claim needs.
\end{intuition}

\subsection{A trend is not a stress effect: the within-year gap}

Correlations between stocks rose over 1990--2022, and stress dates cluster in
a few years (2008, 2020). If coupling simply trends up and stress dates happen
to be late, the stress-minus-calm gap picks up the trend. The fix is a
regression with year fixed effects:
\[
C_t=\theta\,\mathbf 1\{t\in\text{stress}\}+\sum_y\delta_y\mathbf 1\{t\in\text{year }y\}+u_t .
\]
By the Frisch--Waugh--Lovell theorem, $\hat\theta$ is the slope of the
regression of the within-year demeaned $C_t$ on the within-year demeaned
stress dummy: it compares stress and calm dates \emph{of the same year}. The
paper finds that $\hat\theta$ is about \FESHRINK\ smaller than the raw gap, but
still beyond both nulls.

\subsection{Checking size and power by Monte Carlo}

A test has the right \emph{size} if, when the null hypothesis is true, it
rejects in about $5\%$ of samples at the $5\%$ level; its \emph{power} is the
rejection rate when the hypothesis is false. To check both, simulate panels
from known processes, apply the whole test to each, and count rejections:
\begin{itemize}
\item i.i.d.\ Student-$t$, constant correlation: no dynamics at all. Both
nulls should reject about 5\%.
\item GARCH with constant correlation: volatility clustering, no
correlation dynamics. A good test of \emph{dependence} should still reject
about 5\%. The i.i.d.\ null rejects \MCGARCHIID\ of the time, the GARCH null
\MCGARCHFHS.
\item Markov-switching volatility, constant correlation: same question with
regime-type volatility.
\item Markov-switching volatility \emph{and} correlation (0.30 calm, 0.50 in
stress): the effect is real; rejection rates measure power
(\MCPOWFHS\ for the GARCH null).
\end{itemize}

Each null fails somewhere. The i.i.d.\ null over-rejects under GARCH
clustering (17\%); the GARCH null over-rejects when volatility switches
regimes (11\%), because a smooth GARCH recursion cannot copy abrupt
switches. The paper therefore asks a result to pass \emph{both} nulls. Under
this joint rule the worst rejection rate over the three null processes is
6.5\% at the 5\% level and 1.0\% at the 1\% level, while the power against
the correlation-switching process stays at 46\%. With 200 simulations, a
rate of 6.5\% is within about one standard error ($\sqrt{0.05\cdot0.95/200}
\approx1.5$ points) of 5\%.

\subsection{Testing many hypotheses at once: Holm's procedure}

The verdict table tests $m$ statistics. If each is tested at 5\%, the chance
of at least one false rejection grows quickly with $m$. The Holm (1979)
procedure controls this \emph{family-wise error rate}. Sort the $p$-values
$p_{(1)}\le\dots\le p_{(m)}$; reject $H_{(1)}$ if $p_{(1)}\le\alpha/m$, then
$H_{(2)}$ if $p_{(2)}\le\alpha/(m-1)$, and so on, stopping at the first
failure. Equivalently, the adjusted $p$-value of the $k$-th smallest is
$\max_{j\le k}\min\{1,(m-j+1)p_{(j)}\}$. Holm is always at least as powerful
as Bonferroni (which uses $\alpha/m$ for all) and needs no assumption on the
dependence between tests.

\paragraph{A detail specific to resampling $p$-values.} With $B=200$
placebo draws, the smallest possible rank $p$-value is $1/201\approx0.005$.
With $m=23$ tests, Holm requires $p_{(1)}\le0.05/23\approx0.0022$: no rank
$p$-value can ever pass, even for a gap far outside the null distribution.
The paper therefore applies Holm to a Gaussian-tail $p$-value,
$p=1-\Phi(z)$ with $z=(\text{observed}-\text{null mean})/\text{null sd}$,
which can be arbitrarily small. A result passes the verdict table only if
both its rank $p$-value and its Holm-adjusted Gaussian $p$-value are below
5\%. Nine of the 23 tests pass; two results that pass individually (the
compensation slope on industries and the share of episodes where skewness
falls) do not survive the adjustment.

\subsection{Comparing nested forecasts: Clark--West}

H7 compares a small model (volatility only, forecast $f_1$) with a larger one
that adds an entropy regressor (forecast $f_2$). Under the null that the extra
regressor is useless, its true coefficient is zero, but the larger model
\emph{estimates} it, which adds noise: $f_2$ has a slightly larger mean squared
error even when the regressor is irrelevant. The Diebold--Mariano test on
$d_t=(y_t-f_{1,t})^2-(y_t-f_{2,t})^2$ is therefore biased against the larger
model and under-rejects. Clark and West (2007) add back the noise term:
\[
\tilde d_t=(y_t-f_{1,t})^2-\bigl[(y_t-f_{2,t})^2-(f_{1,t}-f_{2,t})^2\bigr],
\]
and test $\E[\tilde d_t]=0$ against $\E[\tilde d_t]>0$ with a HAC standard
error. Why this works: if the extra regressor is useless, $f_2$ equals $f_1$
plus pure estimation noise $\nu_t$ uncorrelated with the forecast error, so
$\E[(y-f_2)^2]=\E[(y-f_1)^2]+\E[\nu^2]$. The DM differential then has mean
$-\E[\nu^2]<0$, and $(f_1-f_2)^2=\nu^2$ is exactly the term that re-centres it
at zero. The test is one-sided and approximately normal.

\begin{papier}
Verdicts (tables 4 and 5 of the journal version): coupling under stress
passes both nulls on both panels, also within years; crisis composition (H8),
the range and daily half-life of $\JJ$ and the largest episode pass the
i.i.d.\ null. One-sided jumps, the shape deficit under stress, the
tails--correlation link (H6) and the daily sign (H5) are produced just as
strongly, or more, by the null.
\end{papier}

%======================================================================
\section{The geometry of stress scenarios}
\label{sec:scenarios}

\subsection{The divergence between two arbitrary Gaussians}

\begin{proposition}
\label{prop:klgauss}
For $p_1=\N(\mu_1,\Sigma_1)$, $p_0=\N(\mu_0,\Sigma_0)$ and $\delta=\mu_1-\mu_0$,
\[
\KL{p_1}{p_0}=\tfrac12\Bigl[\underbrace{\delta^\top\Sigma_0^{-1}\delta}_{\text{shift}}
+\underbrace{\tr(\Sigma_0^{-1}\Sigma_1)-N-\log\det(\Sigma_0^{-1}\Sigma_1)}_{\text{deformation}}\Bigr].
\]
\end{proposition}

\begin{proof}
$\KL{p_1}{p_0}=\E_{p_1}[\log p_1]-\E_{p_1}[\log p_0]$. The first term is
$-h(p_1)=-\tfrac12[N\log2\pi+\log\det\Sigma_1+N]$. For the second, the trace
trick gives
$\E_{p_1}[(X-\mu_0)^\top\Sigma_0^{-1}(X-\mu_0)]=\tr(\Sigma_0^{-1}\Sigma_1)
+\delta^\top\Sigma_0^{-1}\delta$, so
$-\E_{p_1}[\log p_0]=\tfrac12[N\log2\pi+\log\det\Sigma_0+\tr(\Sigma_0^{-1}\Sigma_1)
+\delta^\top\Sigma_0^{-1}\delta]$. Add them and use
$\log\det\Sigma_0-\log\det\Sigma_1=-\log\det(\Sigma_0^{-1}\Sigma_1)$.
\end{proof}

\paragraph{Spectral reading of the deformation.} If $\nu_1,\dots,\nu_N$ are
the (positive) eigenvalues of $\Sigma_0^{-1}\Sigma_1$,
\[
\text{deformation}=\tfrac12\sum_k(\nu_k-1-\log\nu_k).
\]
Each term is $\ge0$ (since $\log\nu\le\nu-1$) and zero if and only if
$\nu_k=1$: every deformed eigendirection costs something.

\paragraph{In one dimension.} For a portfolio with mean $m_0$ and standard
deviation $s_0$ under $p_0$, stressed to $(m,s)$, let $u=(m-m_0)/s_0$ (shift
in standard deviations) and $v=s/s_0$ (volatility multiplier). Then
\begin{equation}
\label{eq:g}
g(u,v)=\tfrac12\bigl(u^2+v^2-1-2\log v\bigr).
\end{equation}

\subsection{The entropy-ball scenario}

\begin{definition}
Given the current law $p_0$, a risk measure $\rho$ and a radius $\eta$, the
scenario of severity $\eta$ is
$p^\star_\eta=\arg\max\{\rho(p):\ \KL{p}{p_0}\le\eta\}$: the worst law within
$\eta$ nats of the current one.
\end{definition}

For a Gaussian loss, VaR and expected shortfall have the form $\rho=-m+c\,s$.
For the $99\,\%$ ES, $c=\varphi(z_{0.99})/0.01=\varphi(2.326)/0.01\approx2.665$.

\begin{proposition}[Solution]
\label{prop:boule}
The maximum of $-m+cs$ subject to $g(u,v)\le\eta$ is reached on the boundary
$g=\eta$, at a point such that
\[
u=-t,\qquad v-\frac1v=c\,t,\qquad t>0 .
\]
For $c=0$ (purely directional shock): $v=1$ and $u=-\sqrt{2\eta}$.
\end{proposition}

\begin{proof}
\emph{Step 1: convexity.} $g$ is the sum of $u^2/2$ (convex) and
$(v^2-1-2\log v)/2$, whose second derivative $1+1/v^2$ is positive: $g$ is
convex on $\R\times(0,\infty)$. The objective $s_0(-u+cv)$ (up to a constant)
is linear. We maximise a linear function over a convex set: the
Karush--Kuhn--Tucker conditions are necessary and sufficient.

\emph{Step 2: the constraint binds.} A non-constant linear function has no
interior maximum; the optimum lies on $g=\eta$.

\emph{Step 3: the Lagrangian.} $\mathcal L=-u+cv-\lambda(g(u,v)-\eta)$. The
first-order conditions are $-1=\lambda u$ and $c=\lambda(v-1/v)$. With
$t=1/\lambda>0$: $u=-t$ and $v-1/v=ct$. The scalar $t$ is pinned down by
$g(-t,v(t))=\eta$.

\emph{Step 4: case $c=0$.} $v-1/v=0$ and $v>0$ force $v=1$, then
$g=u^2/2=\eta$ gives $u=-\sqrt{2\eta}$.
\end{proof}

\begin{exemple}[Checking a row of the severity ladder (online appendix)]
Ten-year rung: $u=-0.38$, $v=1.63$. We check
$g=\tfrac12(0.144+2.657-1-0.977)=0.412\approx\eta=0.408$, and
$v-1/v=1.017\approx c\,t=2.665\times0.38=1.013$. With annual volatility of
$20\,\%$, that is $s_0=0.2/\sqrt{252}=1.26\,\%$ per day, ES goes from
$c\,s_0=3.4\,\%$ to $(0.38+2.665\times1.63)\,s_0=5.9\,\%$. The scenario spends
its budget mostly on volatility ($v$) rather than direction ($u$), because ES
is more sensitive to scale ($c\approx2.7$) than to location (coefficient $1$).
\end{exemple}

\subsection{Calibrating the radius}

We measure how much the market itself revises its law over $h=21$ days:
$I^{(21)}_t=\KL{p_t^w}{p_{t-21}^w}=g(u_t,v_t)$, computed on the portfolio,
with $u_t$, $v_t$ the shift and volatility multiplier observed between $t-21$
and $t$. The history is cut into non-overlapping 21-day blocks (about 12 per
year); a ``ten-year'' radius is the quantile exceeded on average once every
$120$ blocks, that is the $1-1/120$ quantile. Limit: compared with the history
of volatility increases, this radius is about twice too mild; only the
\emph{relative} comparisons that follow are robust.

\subsection{From the portfolio to all assets}

\begin{proposition}[Data-processing inequality]
\label{prop:dpi}
For any function $f$, if $p^f$ and $p_0^f$ are the laws of $f(X)$ under $p$
and $p_0$, then $\KL{p^f}{p_0^f}\le\KL{p}{p_0}$.
\end{proposition}

\begin{proof}[Proof sketch]
Let $Y=f(X)$. The chain rule gives
$\KL{p_{X,Y}}{p_{0,X,Y}}=\KL{p^f}{p_0^f}+\E_{Y}\bigl[\KL{p_{X\mid Y}}{p_{0,X\mid Y}}\bigr]$.
The second term is $\ge0$. Applying the same rule in the other order,
$\KL{p_{X,Y}}{p_{0,X,Y}}=\KL{p}{p_0}+\E_X\bigl[\KL{p_{Y\mid X}}{p_{0,Y\mid X}}\bigr]$,
and this last term is zero because, $Y=f(X)$ being deterministic given $X$,
the two conditional laws coincide. Hence $\KL{p}{p_0}\ge\KL{p^f}{p_0^f}$.
\end{proof}

\begin{intuition}
By looking at less (only the portfolio $w^\top X$), you cannot see more
difference. The cost of a scenario on all assets is always at least its cost
``seen by the portfolio''. Table~7 of the journal version confirms it: $9.97$ nats on 20
assets against $7.24$ seen by the portfolio.
\end{intuition}

\begin{proposition}[Minimal lift]
\label{prop:lift}
Let $p_0=\N(\mu_0,\Sigma_0)$, $s_0^2=w^\top\Sigma_0w$ and
$\beta=\Sigma_0w/s_0^2$ (the vector of asset betas on the portfolio). Among
Gaussians that give the portfolio mean $m_0+us_0$ and standard deviation
$vs_0$, the cheapest is
\[
\mu_1=\mu_0+u\,s_0\,\beta,\qquad\Sigma_1=\Sigma_0+(v^2-1)\,s_0^2\,\beta\beta^\top,
\]
and its cost is exactly $g(u,v)$.
\end{proposition}

\begin{proof}
By proposition~\ref{prop:dpi}, every candidate costs at least $g(u,v)$. It is
enough to show that this candidate attains the bound. We have $w^\top\beta=1$
and $\beta^\top\Sigma_0^{-1}\beta=w^\top\Sigma_0w/s_0^4=1/s_0^2$. Hence:
$\delta^\top\Sigma_0^{-1}\delta=u^2s_0^2\cdot(1/s_0^2)=u^2$;
$\tr(\Sigma_0^{-1}\Sigma_1)=N+(v^2-1)s_0^2\,\beta^\top\Sigma_0^{-1}\beta=N+v^2-1$;
and, by the matrix determinant lemma ($\det(I+ab^\top)=1+b^\top a$),
$\det(\Sigma_0^{-1}\Sigma_1)=1+(v^2-1)s_0^2\,\beta^\top\Sigma_0^{-1}\beta=v^2$.
In proposition~\ref{prop:klgauss}: $\tfrac12(u^2+v^2-1-2\log v)=g(u,v)$.
\end{proof}

\paragraph{Reading.} Write $X=\beta f+e$ with $f=w^\top X$ the portfolio
factor and $e$ uncorrelated with $f$. The minimal scenario touches \emph{only}
the factor: it shifts its mean by $u$ standard deviations and multiplies its
volatility by $v$, leaving $e$ intact. High-beta assets take the shock, and
their correlations rise. \textbf{The least improbable scenario is a
one-factor scenario.}

\begin{proposition}[Cost of uniform dilation]
\label{prop:uniforme}
If all volatilities are multiplied by $v$ with correlations unchanged
($\Sigma_1=v^2\Sigma_0$), the deformation is $N\cdot\tfrac12(v^2-1-2\log v)$:
$N$ times that of the minimal lift.
\end{proposition}

\begin{proof}
Every eigenvalue of $\Sigma_0^{-1}\Sigma_1=v^2I$ equals $v^2$; we sum $N$
identical terms $\tfrac12(v^2-1-\log v^2)$.
\end{proof}

\noindent Numerically, for $v=1.63$ and $N=20$: $0.34$ nat per direction, so
$6.80$ nats in total, against $0.34$ for the minimal lift. This is the
``default: uniform dilation'' row of the composition example in section~6.5 of the journal version.

\subsection{Composing a scenario: volatility or correlation?}

\begin{proposition}
If $w^\top\Sigma_aw=w^\top\Sigma_bw=s^2$, then for every $\alpha\in[0,1]$,
$\Sigma_\alpha=(1-\alpha)\Sigma_a+\alpha\Sigma_b$ satisfies
$w^\top\Sigma_\alpha w=s^2$.
\end{proposition}

\begin{proof}
$\Sigma\mapsto w^\top\Sigma w$ is linear.
\end{proof}

One can therefore mix an ``all-volatility'' covariance and an
``all-correlation'' covariance without changing portfolio severity; $\alpha$
is the \emph{dependence share}. Since $\Sigma_1\mapsto\KL{\N(\mu,\Sigma_1)}{p_0}$
is convex (linear trace, convex $-\log\det$), the cost of a mixture is at most
the mixture of the costs.

\paragraph{A ceiling for correlation.} With volatilities fixed,
$w^\top\Sigma w=\sum_{i,j}w_iw_j\sigma_i\sigma_j\rho_{ij}$ is largest when
every $\rho_{ij}$ equals $1$, and then equals $(\sum_iw_i\sigma_i)^2$. On the
main panel, this is only about $1.5$ times current volatility: beyond that,
volatilities must also be dilated.

\paragraph{Calibrating $\alpha$.} H1 gives the share of dependence in the move
of Gaussian entropy under stress: $0.99/(0.99+5.97)\approx14\,\%$. H8 refines
it to $49\,\%$ over the most severe episodes. The larger $\alpha$, the closer
we get to the one-factor scenario, and the lower the cost
($6.80\to2.88\to0.71$ nats): proposition~\ref{prop:uniforme} makes this result
expected rather than surprising.

\subsection{Pricing a published scenario: the EBA example}

A published scenario sets shocks on some variables $S$ only (equities, long
rates), not on the others $U$. The minimal cost is obtained by letting the
unshocked variables take their most likely value.

\begin{proposition}[Cost of a partial shock]
\label{prop:partiel}
For $p_0=\N(0,\Sigma)$ and a shift $\delta=(\delta_S,\delta_U)$ of which only
$\delta_S$ is imposed,
\[
\min_{\delta_U}\ \tfrac12\delta^\top\Sigma^{-1}\delta
=\tfrac12\,\delta_S^\top\Sigma_{SS}^{-1}\delta_S ,
\qquad\text{reached at }\delta_U^\star=\Sigma_{US}\Sigma_{SS}^{-1}\delta_S .
\]
\end{proposition}

\begin{proof}
Let $P=\Sigma^{-1}$, partitioned into blocks. The function
$\delta_U\mapsto\delta^\top P\delta$ is strictly convex quadratic; its
derivative vanishes at $\delta_U^\star=-P_{UU}^{-1}P_{US}\delta_S$, and the
minimum is $\delta_S^\top(P_{SS}-P_{SU}P_{UU}^{-1}P_{US})\delta_S$. The matrix
in brackets is the Schur complement of $P_{UU}$ in $P$, and the block
inversion formula says precisely that it equals $\Sigma_{SS}^{-1}$. Finally
$-P_{UU}^{-1}P_{US}=\Sigma_{US}\Sigma_{SS}^{-1}$ (same formula): this is the
Gaussian conditional expectation of $X_U$ given $X_S=\delta_S$.
\end{proof}

\begin{exemple}[Why correlation makes a shock cheaper]
Two indices with the same volatility, correlated at $\rho$, each fall by $k$
standard deviations. With $\delta=(k,k)$ and $R^{-1}=\frac1{1-\rho^2}
\begin{pmatrix}1&-\rho\\-\rho&1\end{pmatrix}$,
\[
\tfrac12\delta^\top R^{-1}\delta=\frac{k^2}{1+\rho},
\]
whereas the sum of the costs taken one by one is $k^2$. With $\rho=0.87$
(average correlation of European indices), the joint cost is only $53\,\%$ of
the sum. This is the mechanism that takes the EBA scenario from $31.9$ nats
(variable by variable) to $10.5$ nats (joint): a common shock to highly
correlated indices follows the most likely direction in the data.
\end{exemple}

\paragraph{The EBA figures.} For the DAX alone, $\Delta=\log0.5=-0.693$ and
$\sigma=20.3\,\%$ per year: the shock is $3.41$ standard deviations and costs
$\tfrac12(3.41)^2=5.8$ nats. The full scenario costs $10.5$ nats, that is
$\sqrt{2\times10.5}=4.6$ standard deviations on a single factor.

\begin{exercice}
Show that the cost of a $k$-standard-deviation fall in $N$ indices with
constant correlation $\rho$ is $\tfrac{N k^2}{2(1+(N-1)\rho)}$.
\end{exercice}
\noindent\emph{Solution.} $\delta=k\mathbf 1$ is an eigenvector of $R$ with
eigenvalue $1+(N-1)\rho$, so $\delta^\top R^{-1}\delta=k^2N/(1+(N-1)\rho)$.

\subsection{Reverse stress testing}

We reverse the question: what is the cheapest scenario that produces a given
loss $L$? In the Gaussian setting, it is proposition~\ref{prop:boule} read
backwards. In the non-parametric setting, the solution is an
\emph{exponential tilting} of the sample: observed days are reweighted with
weights proportional to $e^{\theta\,\text{loss}_t}$. As $\theta$ grows, almost
all the weight goes to a few days; the effective sample size
$(\sum_t\omega_t)^2/\sum_t\omega_t^2$ collapses (it is $1.0$ at $9$ nats on
the main panel), and one must go back to the Gaussian version.

%======================================================================
\section{Decomposition by segment: the entropic CISS}
\label{sec:blocs}

Split the $N$ series into blocks $B_1,\dots,B_K$ (money market, bonds,
equities, intermediaries, FX). Let $R_{kk}$ be the diagonal blocks of $R$ and
$R_D=\operatorname{diag}(R_{11},\dots,R_{KK})$ the block-diagonal matrix.

\begin{proposition}
\label{prop:blocs}
\[
\underbrace{C(R)}_{\text{total}}
=\sum_{k=1}^K\underbrace{C(R_{kk})}_{\text{within block }k}
+\underbrace{\Bigl(-\tfrac12\log\det\bigl(R_D^{-1/2}RR_D^{-1/2}\bigr)\Bigr)}_{\text{between blocks}},
\]
and each of the $K+1$ terms is non-negative. The between-block term is the
Gaussian mutual information between the blocks; it is zero if and only if the
blocks are pairwise uncorrelated.
\end{proposition}

\begin{proof}
\emph{Identity.} $R=R_D^{1/2}\,\tilde R\,R_D^{1/2}$ with
$\tilde R=R_D^{-1/2}RR_D^{-1/2}$. Since the determinant is multiplicative,
$\log\det R=\log\det R_D+\log\det\tilde R$, and
$\log\det R_D=\sum_k\log\det R_{kk}$ (determinant of a block-diagonal matrix).
Multiply by $-\tfrac12$.

\emph{Signs.} Each $R_{kk}$ is a correlation matrix, so $C(R_{kk})\ge0$ by
proposition~\ref{prop:hadamard}. The matrix $\tilde R$ has identity diagonal
blocks, hence a diagonal of $1$s: it is again a correlation matrix, and the
same argument applies. (Equivalently, $\det R\le\prod_k\det R_{kk}$ is
Fischer's inequality.)
\end{proof}

\paragraph{Two blocks: canonical correlations.} For $K=2$, let
$M=R_{11}^{-1/2}R_{12}R_{22}^{-1/2}$; its singular values $\rho_1,\dots,\rho_r$
are the \emph{canonical correlations} between the two segments. We have
$\tilde R=\begin{pmatrix}I&M\\M^\top&I\end{pmatrix}$ and, by the Schur
complement formula, $\det\tilde R=\det(I-MM^\top)=\prod_j(1-\rho_j^2)$. Hence
\[
I(B_1;B_2)=-\tfrac12\sum_j\log(1-\rho_j^2):
\]
the sum of the information carried by each canonical pair, each computed as
in proposition~\ref{prop:migauss}.

\begin{intuition}
Total coupling splits exactly into ``what happens inside each market'' and
``what ties markets together''. The second term is the systemic part. Unlike
the CISS, which aggregates, this decomposition says \emph{which} segments
stress ties together, and by how much, in a single unit.
\end{intuition}

\paragraph{The rank version.} Replace each series by
$\Phi^{-1}(\hat F_i(x_i))$, where $\hat F_i$ is its empirical distribution
function: this gives ``normal scores''. A monotone transformation of the
marginals does not change ranks, so the measure becomes insensitive to each
series' fat tails: it estimates the multi-information of a \emph{Gaussian
copula}. This is the most robust version against the null ($9.6$ and $12.7$
standard deviations).

\paragraph{Link with the structural index.} With the chain rule,
\[
\JJ\approx\underbrace{\textstyle\sum_i\DD_i}_{\text{shape}}
+\underbrace{\textstyle\sum_kC(R_{kk})}_{\text{within}}
+\underbrace{(-\tfrac12\log\det\tilde R)}_{\text{between}} .
\]
The entropic CISS keeps the last two terms and drops shape, which is what the
null reproduces. It is the structural index stripped of its non-robust
channel.

%======================================================================
\section{Summary: each formula and its role}
\label{sec:recap}

\begin{center}
\small
\begin{tabular}{@{}p{0.42\linewidth}p{0.52\linewidth}@{}}
\toprule
Result & Role in the paper \\
\midrule
$h(aX)=h(X)+\log|a|$ & the level of entropy has no absolute meaning; we work
with divergences \\
$\KL{p}{q}\ge0$, invariant under change of variables & common unit across all
asset classes \\
$\tfrac12\Delta^\top\Sigma^{-1}\Delta$; $k\sigma\leftrightarrow k^2/2$ nats &
nats--standard deviations exchange rate; scenario audit \\
Gaussian maximum entropy & the Gaussian is the zero point \\
$h(q)=\tfrac N2\log(2\pi e)+\sum\log\sigma_i+\tfrac12\log\det R$ & separating
scale and dependence \\
$\det R\le1$ & coupling $C$ is non-negative \\
$\DD=h(q)-h(p)$; $\JJ=\DD+C=\KL{p}{\prod q_i}$ & the structural index is a
true divergence \\
$\E_Y[h(X\mid Y)]=h(X)-I(X;Y)$ & upward jumps of $\JJ$ (theorem) \\
OU, $t_{1/2}=\ln2/\kappa$ & relaxation (postulate), half-life \\
Jensen on a convex estimator & upward bias; need for a null \\
Bartlett: $\E[C(\hat R)]\approx N(N-1)/(4T')$ & size of the bias; compare gaps, not levels \\
Matrix determinant lemma, one common day & the mechanical jump of coupling \\
Exchangeability, $(1+k)/(B+1)$ & exact resampling $p$-values \\
CCC-GARCH filtered historical simulation & null that keeps volatility, not correlation, dynamics \\
Year fixed effects (Frisch--Waugh--Lovell) & stress effect within a year, not a trend \\
Holm step-down & family-wise error control in the verdict table \\
Clark--West & forecast test for nested models (H7) \\
KL of two Gaussians, $g(u,v)$ & cost of a scenario on a portfolio \\
KKT: $u=-t$, $v-1/v=ct$ & entropy-ball scenario \\
Data processing; minimal lift & the minimal scenario is one-factor \\
Uniform dilation $=N\times$ minimum & why classical scenarios are too
expensive \\
Schur complement & cost of a partial shock (EBA) \\
Within / between-block decomposition & entropic CISS, contagion map \\
\bottomrule
\end{tabular}
\end{center}

\paragraph{Further reading.} Cover and Thomas, \emph{Elements of Information
Theory} (2006), chapters 2, 8 and 12, cover entropy, divergence, differential
entropy and maximum entropy. Breuer and Csiszár (2013) and Glasserman and Xu
(2014) develop Kullback--Leibler ball scenarios. The project's theoretical
supplement gives the proofs in the paper's exact notation.

\end{document}
